\section{Choice of parameters and proof of Theorem~\ref{thm:main}}
\label{sec:parameters}

We now combine the interpolation theorem and the determinant estimates.
The order of choices is essential: the dimension is fixed before the
approximation scales, and all these data are fixed before the polynomial
degree tends to infinity.

\begin{lemma}\label{par:constants}
For every real \(\nu>2\), there are positive rational numbers
\(\theta,A,B,C\) such that
\begin{equation}\label{par:inequalities}
0<\theta<A<B<1,\qquad
\nu(A-\theta)>1-\theta,\qquad
C>1,\quad B<1/C,\quad B<C\theta<1.
\end{equation}
One can then choose \(0<\eta<1\) such that
\begin{equation}\label{par:gap}
g:=\nu\bigl(A(1-\eta)-\theta\bigr)-(1-\theta)>0.
\end{equation}
\end{lemma}

\begin{proof}
Choose a rational \(b\) with \(1/2<b<1-1/\nu\).
For a sufficiently small positive rational \(\delta\), put
\(\theta=1-\delta\) and \(A=1-b\delta\). Then \(0<\theta<A<1\), and
\[
\nu(A-\theta)-(1-\theta)
=\bigl(\nu(1-b)-1\bigr)\delta>0,
\]
whereas
\[
\theta-A^2=(2b-1)\delta-b^2\delta^2>0.
\]
It follows that the interval \((A/\theta,1/A)\) is nonempty.
Choose a rational \(C\) in this interval. Since \(A>\theta\), we have
\(C>1\); since \(\theta<A\), we also have
\(C\theta<\theta/A<1\).
Now choose a rational \(B\) between \(A\) and
\(\min(1/C,C\theta)\). This proves \eqref{par:inequalities}.
Finally, the strict inequality at \(\eta=0\) allows a sufficiently
small \(\eta>0\) in \eqref{par:gap}.
\end{proof}

The two inequalities used in this construction have different roles.
The inequality \(\nu(A-\theta)>1-\theta\) supplies the
approximation-error gap. The inequality \(A^2<\theta\) permits
\(A/\theta<C<1/A\), and hence a choice of \(B\) for which
\(CB<1\), \(C\theta/B>1\), and \(B/A>1\). With the full-simplex
counts of Section~\ref{det:section}, these three ratios make the
dimension-dependent errors tend to zero while the collision saving
tends to infinity. We now choose one finite dimension realizing
both requirements; it will remain fixed throughout the degree limit.

\begin{proof}[Proof of Theorem~\ref{thm:main}]
Fix \(\nu>2\). Suppose for a contradiction that there are arbitrarily
large positive integers \(q\) for which some \(p\in\ZZ\) satisfies
\begin{equation}\label{par:good}
\left|\pi-\frac pq\right|\le q^{-\nu}.
\end{equation}
Fix \(\theta,A,B,C,\eta,g\) from Lemma~\ref{par:constants}, and put
\(\varepsilon_0=\min(g/2,1/2)\).
First choose \(F_0>2/\theta\) so large that
\(\nu/F_0<\varepsilon_0/3\).

For an integer \(m\), set
\begin{equation}\label{par:dimension}
K=\lfloor C^m\rfloor,\qquad
w_0=B^{-m},\qquad
v_0=2K\theta^m w_0.
\end{equation}
These are admissible rational interpolation data:
\begin{equation}\label{par:volume}
K\theta^m\le(C\theta)^m<1,\qquad
K\frac{w_0}{v_0}\theta^m=\frac12.
\end{equation}
The relevant limits as \(m\to\infty\) are
\begin{equation}\label{par:limits}
\frac{K}{w_0}\le(CB)^m\longrightarrow0,\qquad
v_0\sim2\left(\frac{C\theta}{B}\right)^m\longrightarrow\infty,
\qquad
\frac{m}{v_0}\longrightarrow0.
\end{equation}
With the notation of Proposition~\ref{det:comparison}, these choices give
\begin{equation}\label{par:collision-growth}
L=\frac{\eta^2K\theta^m}{(m+1)v_0A^m}
 =\frac{\eta^2(B/A)^m}{2(m+1)}\longrightarrow\infty.
\end{equation}
Thus, with \(\Lambda=4\log2\) and \(c=(\log2)/4\), choose \(m\) large
enough that
\begin{equation}\label{par:dimension-margin}
\frac{\Lambda F_0m+2\log2}{v_0}+\frac{100K}{w_0}
 <\frac{\varepsilon_0}{3},
\qquad cL>2.
\end{equation}
Fix this \(m\), and hence \(K,w_0,v_0\).

The remaining constants are now fixed. Choose approximations
\(p_i/q_i\) satisfying \eqref{par:good}, with
\(w_i=\lceil\log q_i\rceil\), successively sufficiently large for
Theorem~\ref{geom:interpolation} and, in addition, for
\begin{equation}\label{par:weight-margin}
\Lambda\sum_{i=1}^m\frac1{w_i}
+\frac{\theta+\log4+\log(2K)+\nu+\log(200K)}{w_*}
 <\frac{\varepsilon_0}{3},
\qquad w_*=\min_{i>0}w_i.
\end{equation}
This is possible: first impose a sufficiently large lower bound on
every \(w_i\), using \(\sum_i1/w_i\le m/w_*\), and then impose the
successive geometric thresholds. The assumed denominators are
unbounded, so every finite sequence of such requirements can be met.
The thresholds for weight separation are independent of the centers.
For sufficiently large \(q_i\), \(p_i\ne0\), and hence
\(c_{ji}=jr_i\), \(r_i=2\ii p_i/q_i\), are distinct across \(j\) in
each coordinate. No independence between different coordinates is
required.

Using \eqref{det:arith-error}, \eqref{det:translation-bound},
\eqref{det:hol-error}, and \eqref{det:analytic-error}, the sum of errors
is exactly
\[
\begin{aligned}
E_{\rm ar}+E_{\rm an}
&=\frac{\nu}{F_0}
 +\frac{\Lambda F_0m+2\log2}{v_0}+\frac{100K}{w_0}\\
&\quad +\Lambda\sum_i\frac1{w_i}
 +\frac{\theta+\log4+\log(2K)+\nu+\log(200K)}{w_*}.
\end{aligned}
\]
Our three choices make this less than \(\varepsilon_0\).
In particular it is less than \(g\), while
\(cL>2>1+E_{\rm ar}+E_{\rm an}\).
All hypotheses of Proposition~\ref{det:comparison} now hold,
a contradiction. The degree \(H\) in that Proposition tends to infinity
only after the dimension, weights, and centers just chosen are fixed.

We have excluded \eqref{par:good} for arbitrarily large \(q\), for
each fixed \(\nu>2\). This gives the required eventual lower bound,
in fact with a strict inequality. The contradiction allowed unreduced
fractions and zero error. If \(\pi\) were rational, exact representations
with arbitrarily large denominators would contradict it; hence
\(\pi\) is irrational.

Finally divide the unit interval into \(N\) equal parts and apply
pigeonhole to \(0,\pi,\ldots,N\pi\) modulo one. There exist \(p\in\ZZ\)
and \(1\le q\le N\) such that
\[
0<\left|\pi-\frac pq\right|\le\frac1{qN}\le\frac1{q^2}.
\]
After reduction, the same error is at most the inverse square of the
reduced denominator. These errors tend to zero as \(N\to\infty\);
irrationality forces the reduced denominators to be unbounded.
Every exponent less than two therefore occurs infinitely often in
the defining inequality. Hence \(\mu(\pi)\ge2\), and the upper bound
already proved gives \(\mu(\pi)=2\).
\end{proof}
